\documentclass[11pt]{report}
\usepackage[T5]{fontenc}
\usepackage{helvet}
\usepackage{caption}
\usepackage{amssymb}
\usepackage[utf8]{inputenc}
\usepackage[vietnamese,english]{babel}
\usepackage{amsmath}
\usepackage{graphicx}
\usepackage[colorinlistoftodos]{todonotes}

\usepackage{changepage}

\usepackage{float}

\usepackage{lipsum} 
\usepackage[]{mdframed}

\usepackage[shortlabels]{enumitem}

\usepackage{listings}
\usepackage{caption}

\usepackage{pgfplots}
\pgfplotsset{width=10cm,compat=1.9}

\usepackage{parskip}
\setlength{\parskip}{0pt}

\usepackage{multirow}
\usepackage{multicol}

\captionsetup{justification=centering}
\usepackage{setspace}
\usepackage{hyperref}
\usepackage{float}
\usepackage{amsmath}
\usepackage{titlesec}
\hypersetup{
    colorlinks=true,
    linkcolor=blue,
    filecolor=magenta,      
    urlcolor=cyan,
}

\usepackage{adjustbox}
\usepackage{xcolor}
\setlength{\parskip}{1em}  % Adjusts the space between paragraphs to 1em
\definecolor{darkgreen}{rgb}{0,0.39,0}
\lstset{
  basicstyle=\itshape,
  columns=fullflexible,
  xleftmargin=3em,
  breaklines=true,
  postbreak=\mbox{\textcolor{red}{$\hookrightarrow$}\space},
  mathescape,
  literate={->}{$\rightarrow$}{2}
}
\usepackage{indentfirst} % Đảm bảo rằng đoạn đầu tiên sau tiêu đề cũng được thụt đầu dòng
\setlength{\parindent}{0.5cm}

\usepackage{fancyhdr}
\usepackage{geometry}
\usepackage{enumitem}

% Define page geometry
\geometry{top=1in, bottom=1in, left=1in, right=1in, headheight=15pt}
\geometry{footskip=30pt}

% % Define subsubsubsection
% \titleclass{\subsubsubsection}{straight}[\subsubsection]
% \newcounter{subsubsubsection}[subsubsection]
% \renewcommand{\thesubsubsubsection}{\thesubsubsection.\arabic{subsubsubsection}}
% \titleformat{\subsubsubsection}
%   {\normalfont\normalsize\bfseries}{\thesubsubsubsection}{1em}{}
% \titlespacing{\subsubsubsection}{0pt}{3.25ex plus 1ex minus .2ex}{1em}

%section numbering
% Custom formatting for sections, subsections, etc.
\setcounter{secnumdepth}{3}
\renewcommand{\thesection}{\arabic{section}}
\renewcommand{\thesubsection}{\thesection.\arabic{subsection}}
\renewcommand{\thesubsubsection}{\thesubsection.\arabic{subsubsection}}

% Define and manage custom counters for subsections and subsubsections
\newcounter{subsectioncount}
\newcounter{subsubsectioncount}

\let\oldsubsection\subsection
\let\oldsubsubsection\subsubsection

\renewcommand{\subsection}[1]{%
  \stepcounter{subsectioncount}%
  \setcounter{subsubsectioncount}{0}%
  \oldsubsection{#1}%
}

\renewcommand{\subsubsection}[1]{%
  \stepcounter{subsubsectioncount}%
  \oldsubsubsection{#1}%
}


%to include subsubsection in table of contents
\setcounter{tocdepth}{3}

% Setup fancy headers and footers
\pagestyle{fancy}
\fancyhf{} % clear all header and footer fields
\renewcommand{\headrulewidth}{0.4pt} % line in header area
\renewcommand{\footrulewidth}{0.4pt} % line in footer area
\fancyhead[L]{CSC14119} % Header left
\fancyhead[R]{Course Milestone 1 Report} % Header right
%\fancyfoot[C]{\uppercase{\rightmark}} % Footer center
\fancyfoot[R]{\thepage} % Footer right

%can giua TOC
\usepackage{tocloft}

% Center the TOC title
\renewcommand{\cfttoctitlefont}{\hfill\Large\bfseries}
\renewcommand{\cftaftertoctitle}{\hfill}

% Adjust TOC entry layout if needed
%\setlength{\cftsecindent}{0pt} % Set section indentation to 0
%\setlength{\cftsubsecindent}{5pt} % Set subsection indentation to 0
%\setlength{\cftsubsubsecindent}{10pt} % Set subsubsection indentation to 0

\begin{document}

\begin{titlepage}
\textrm{}
\newcommand{\HRule}{\rule{\linewidth}{0.3mm}} % Defines a new command for the horizontal lines, change thickness here

\center % Center everything on the page
 

\textsc{\LARGE UNIVERSITY OF SCIENCE}\\[1.0cm] % Name of your university/college

\textsc{\Large Faculty of Information Technology}\\[0.5cm] % Major heading such as course name

\textsc{\large \bfseries CSC14119 - Introduction to Data Science}\\[1.5cm] % Minor heading such as course title

%----------------------------------------------------------------------------------------
%	TITLE SECTION
%----------------------------------------------------------------------------------------

\HRule \\[0.2cm]
\LARGE \bfseries REPORT\\[0.4cm]
{ \fontsize{30pt}{30pt}\selectfont \bfseries COURSE MILESTONE 1}\\[0.4pt]
\HRule \\[1.0cm]
 
%----------------------------------------------------------------------------------------
%	AUTHOR SECTION
%----------------------------------------------------------------------------------------

\nopagebreak
\large 23KHDL2

\begin{minipage}{0.4\textwidth}
\begin{flushleft} \large
\emph{Students:}\\
Nguyễn Đỗ Bảo - 23127026 \\
Vũ Hoàng Minh - 23127427 \\
Nguyễn Thanh Owen - 23127447
\end{flushleft}
\end{minipage}
~
\begin{minipage}{0.4\textwidth}
\begin{flushright} \large
\emph{Instructors:} \\
Mr. Huỳnh Lâm Hải Đăng % Supervisor's Name
\end{flushright}
\end{minipage}\\[2cm]


 

\vfill % Fill the rest of the page with whitespace

\end{titlepage}

\newpage 
\pagestyle{fancy}
\fancypagestyle{plain}{%
\fancyhf{} % clear all header and footer fields
\renewcommand{\headrulewidth}{0.4pt} % line in header area
\renewcommand{\footrulewidth}{0.4pt} % line in footer area
\fancyhead[L]{CSC14119} % Header left
\fancyhead[R]{Course Milestone 1 Report} % Header right
%\fancyfoot[C]{\uppercase{\rightmark}} % Footer center
\fancyfoot[R]{\thepage} % Footer right
}
\tableofcontents
%\thispagestyle{fancy}
\vfill
\newpage
%\renewcommand{\thesection}{\arabic{section}} % Set section numbering format

%section 1
\section{Overview}
This scraper is designed to systematically collect arXiv papers, their source files, metadata, and references within a specified ID range. The implementation follows a multi-threaded approach to optimize performance while respecting API rate limits.

%section 2
\section{Tools and Libraries Used}
\begin{itemize}
    \item \texttt{arxiv}: Python wrapper for arXiv API to query and download papers
    \item \texttt{requests}: HTTP library for Semantic Scholar API calls
    \item \texttt{psutil}: System and process utilities for memory monitoring
    \item \texttt{concurrent.futures} and \texttt{threading}: Multi-threading for parallel paper processing
    \item \texttt{tarfile}: Built-in library for extracting .tar.gz source archives
    \item \texttt{json}: Built-in library for metadata and reference storage
    \item \texttt{urllib.error}: Handle URL and HTTP response errors.
\end{itemize}

\section{Scraping Workflow}
\begin{enumerate}
    \item Entry Discovery: Uses arXiv API to search for papers by ID within the specified range (2305.14597 to 2306.09504). Binary search is employed to find the last existing paper ID in each month.
    \item Source Download: For each paper, all versions of which are downloaded using the arXiv API's \texttt{download\textunderscore source()} method by patching the PDF URL for each version number.
    \item Extraction and Filtering: Source archives (\texttt{.tar.gz}) are extracted, retaining only \texttt{.tex} and \texttt{.bib} files while removing figures and other non-essential files to reduce storage.
    \item Metadata Collection: Paper metadata (title, authors, submission date, revised dates, publication venue) is retrieved via arXiv API with timezone conversion to US Eastern Time for accurate revision dates.
    \item Reference Extraction: Semantic Scholar API is used to retrieve reference lists with arXiv IDs, including retry logic for rate limit handling (429 errors).
\end{enumerate}

\section{Statistics}
This statistics is the result of scraping the papers by ID within the range of 2306.4505 to 2306.6004.
\subsection{General Information}
\begin{table}[H]
\centering
   \begin{tabular}{|l|l|}
        \hline
        \textbf{Criteria} & \textbf{Stats} \\
        \hline
        Papers processed & 1500 \\
        \hline
        Total papers processed successfully & 1500 \\
        \hline
        Total time & 6863.95s \\
        \hline
        Average time per paper & 4.58s \\
        \hline
        Maximum memory usage & 44.76MB \\
        \hline
        Average memory usage & 44.45MB \\
        \hline
        Paper scraping success rate & 100.00\% \\
        \hline
    \end{tabular}   
    \caption{General statistics about the result}
\end{table}

\subsection{Source and Extraction}
\begin{table}[H]
\centering
   \begin{tabular}{|l|l|}
        \hline
        \textbf{Criteria} & \textbf{Stats} \\
        \hline
        Total versions downloaded & 2479 \\
        \hline
        Average versions per paper & 1.65 \\
        \hline
        Total source size  & 10761914296 bytes \\
        \hline
        Average size per paper & 7174609.53 bytes \\
        \hline
        Total size after extraction & 1723361378 bytes \\
        \hline
        Average size after extraction per paper & 1148907.59 bytes\\
        \hline
    \end{tabular}
    \caption{Source and extraction size statistics}
\end{table}

\subsection{Metadata}
\begin{table}[H]
\centering
   \begin{tabular}{|l|l|}
        \hline
        \textbf{Criteria} & \textbf{Stats} \\
        \hline
        Total metadata collected & 1500 \\
        \hline
        Total metadata size  & 596037 bytes \\
        \hline
        Average metadata size per paper & 397.36 bytes \\
        \hline
    \end{tabular}
    \caption{Metadata statistics}
\end{table}

\subsection{References}
\begin{table}[H]
\centering
   \begin{tabular}{|l|l|}
        \hline
        \textbf{Criteria} & \textbf{Stats} \\
        \hline
        Total references collected & 32476 \\
        \hline
        Average references per paper & 21.65 \\
        \hline
        Total references file size & 13688357 bytes \\
        \hline
        Reference scraping success rate & 99.60\% \\
        \hline
    \end{tabular}
    \caption{References statistics}
\end{table}

\subsection{Disk Usage (Extract - Metadata - References)}
\begin{table}[H]
\centering
   \begin{tabular}{|l|l|}
        \hline
        \textbf{Criteria} & \textbf{Stats} \\
        \hline
        Total disk usage & 1737645772 bytes \\
        \hline
        Average disk usage per paper & 1158430.51 bytes \\
        \hline
    \end{tabular}
    \caption{References statistics}
\end{table}

\section{Demo Video}
Video Link: \href{https://youtu.be/XRfLFomCfhw}{Demo Video}

\end{document}